% =====================================================================
%  Stage 2, раздел 3.
%  Синтетические тестовые сценарии и протокол экспериментов.
% =====================================================================

\section{Синтетические тестовые сценарии и протокол экспериментов}

\subsection{Назначение раздела}

В разделе формализованы (а) библиотека синтетических тестовых сценариев,
формируемая ИИ-злоумышленником, и (б) протокол экспериментов Этапа~4 ВКР.
Раздел является связующим звеном между методическими решениями
разделов~1--2 настоящего отчёта (Этап~2) и реализацией прототипа на
Этапе~3.

\subsection{Библиотека синтетических атакующих сценариев}

В соответствии с обоснованием раздела~3 Этапа~1, библиотека реализуется
в виде набора параметризованных шаблонов, согласованных с таксономией
UNSW-NB15. Перечень и параметры шаблонов приведены в
\tabref{tab:attack-templates}.

\begin{table}[H]
\centering
\caption{Библиотека синтетических атакующих шаблонов}
\label{tab:attack-templates}
\renewcommand{\arraystretch}{1.25}
\setlength{\tabcolsep}{4pt}
\small
\begin{tabularx}{\linewidth}{@{}lXX@{}}
\toprule
\textbf{Шаблон} & \textbf{Параметризация} & \textbf{Соответствие классу UNSW-NB15} \\
\midrule
\texttt{ddos\_synflood} &
интенсивность \(\lambda_0\), длительность \(\Delta\), доля SYN-флагов,
распределение исходных IP &
DoS \\
\addlinespace
\texttt{port\_scan} &
скорость подключений \(\rho\), диапазон портов, режим (последовательный /
случайный), длительность сессии &
Reconnaissance \\
\addlinespace
\texttt{brute\_force} &
интервал между попытками, целевой порт (22, 3389, 80, 443),
доля успешных авторизаций &
Reconnaissance / Exploits \\
\addlinespace
\texttt{exploit\_payload} &
размер payload, нестандартные сочетания TCP-флагов, доля retransmissions &
Exploits / Backdoors \\
\addlinespace
\texttt{data\_exfil} &
длительность сессии, асимметрия \(b_{out}/b_{in}\), стабильность скорости,
прикладной сервис &
Backdoors / Worms \\
\addlinespace
\texttt{internal\_recon} &
фан-аут потоков (число destinations за период), периодичность,
размер сессии &
Reconnaissance \\
\addlinespace
\texttt{worm\_lateral} &
число направленных потоков, периодичность, размер сессии,
распределение destinations &
Worms \\
\bottomrule
\end{tabularx}
\end{table}

Каждый шаблон возвращает упорядоченную последовательность flow-векторов,
семантически согласованную с признаковой схемой UNSW-NB15. Параметры по
умолчанию выбраны из обоснованных в литературе диапазонов
характеристик каждого класса
атак~\cite{moustafa2015unsw, sharafaldin2018toward, sharafaldin2019ddos,
khraisat2019survey}.

\subsection{Конечно-автоматная политика агента}

Конечный автомат агента, формализованный в разделе~3 Этапа~1, в текущем
протоколе принимает следующее множество состояний:
\[
\mathcal{S} =
\{\text{NORMAL}\} \cup
\{\text{ATTACK}(c)\}_{c \in \{\,\text{ddos, scan, brute, exploit, exfil, intrec, worm}\,\}} \cup
\{\text{DRIFT}(k)\}_{k \in \{\,\text{cov, prior, concept}\,\}}.
\]

Политика \(\pi\) задаётся вероятностями переходов, минимальными длительностями
состояний и параметрами рандомизации шаблонов.

\subsection{Протокол валидации генератора}

Перед использованием в экспериментах E1--E6 (см. ниже) генератор проходит
валидацию по двум критериям, формализованным в разделе~3 Этапа~1:

\begin{enumerate}
    \item \textbf{Маргинальное соответствие нормы.} Двувыборочный тест
    Колмогорова--Смирнова и MMD-сравнение распределений признаков
    \emph{dur, sbytes, dbytes, sttl, sintpkt, dintpkt} режима NORMAL агента
    с распределениями нормального трафика обучающей выборки UNSW-NB15.
    Условие приёма: $p\text{-value}_{\text{KS}} > 0{,}01$ и MMD не
    превышает 95-перцентиля bootstrap-распределения.
    \item \textbf{Семантическая согласованность атак.} Сравнение средних и
    медиан ключевых признаков (\emph{dur, sbytes, dbytes, sintpkt}) для
    каждого ATTACK-состояния с соответствующим классом UNSW-NB15
    (DoS, Reconnaissance, Exploits, Backdoors, Worms). Условие приёма:
    относительное отклонение средних не более 25\% и совпадение знака
    асимметрии распределений.
\end{enumerate}

При невыполнении условий приёма параметры шаблона корректируются до
выполнения; результаты валидации фиксируются в журнале и приводятся в
отчёте Этапа~4.

\subsection{Протокол экспериментов}

В \tabref{tab:experiments} зафиксированы шесть экспериментов с указанием
цели, входов, метрик и условий принятия гипотезы.

\begin{table}[H]
\centering
\caption{Протокол экспериментов Этапа~4 ВКР}
\label{tab:experiments}
\renewcommand{\arraystretch}{1.25}
\setlength{\tabcolsep}{4pt}
\small
\begin{tabularx}{\linewidth}{@{}lXXX@{}}
\toprule
\textbf{ID} & \textbf{Цель и условия} & \textbf{Метрики} & \textbf{Критерий принятия} \\
\midrule
E1 &
Сравнение моделей-кандидатов (CNN-LSTM, BiLSTM, LSTM, GRU, 1D-CNN, MLP,
Transformer, AE, VAE, RF, XGBoost, LR, IF, OCSVM) на UNSW-NB15;
\(\geq 5\) seed; одинаковый препроцессинг и пороги &
Precision, Recall, F1, MCC, PR-AUC, ROC-AUC; \(\sigma_{F1}\) &
\(F_1^{\text{CNN-LSTM}} \geq 0{,}90\); статистически значимое
($p < 0{,}05$, парный Уилкоксон) превосходство над ансамблем baselines \\
\addlinespace
E2 &
Анализ ошибок (FP/FN) лучшей модели; разрезы по \emph{attack\_cat},
по протоколу, по сервису &
confusion matrix, FN/FP по классам, ROC по классам &
явное отображение классов, в которых F1 ниже среднего значения по тесту \\
\addlinespace
E3 &
Калибровка порога; температурное масштабирование &
ECE, FP-rate vs Recall, FP-per-time, TTD &
наличие рабочей точки, удовлетворяющей NFR-2 \\
\addlinespace
E4 &
Тестирование с ИИ-злоумышленником на стенде; 6 шаблонных сценариев
+ 3 drift-сценария &
Precision, Recall, F1, TTD, FP-per-time;
ground truth из FSM-журнала &
\(\Delta F_1 \leq 0{,}10\) при agent-NORMAL+ATTACK без drift; TTD$\leq W$
для шаблонов с явными flow-проявлениями \\
\addlinespace
E4' &
Cross-evaluation: train UNSW $\to$ test CICIDS2017
без дообучения и c дообучением на 10\% &
Precision, Recall, F1, PR-AUC, $\Delta F_1$ относительно UNSW &
$\Delta F_1 \leq 0{,}20$ без дообучения; $\leq 0{,}10$ с дообучением \\
\addlinespace
E5 &
Drift-устойчивость; PSI/MMD-инжектирование &
PSI, MMD, $\Delta F_1$ vs PSI; время восстановления (recovery time) &
$\Delta F_1 \leq 0{,}10$ при PSI $\leq 0{,}25$; срабатывание
drift-monitor \\
\addlinespace
E6 &
Эксплуатационная производительность &
inference latency (ms / window),
throughput (EPS), CPU/RAM &
NFR-3, NFR-4 \\
\bottomrule
\end{tabularx}
\end{table}

\subsection{Меры обеспечения статистической достоверности}

В соответствии с замечаниями TESSERACT~\cite{pendlebury2019tesseract} и
обзоров~\cite{khraisat2019survey, sommer2010outside} для каждого эксперимента
реализуются:

\begin{enumerate}
    \item \textbf{Множественные seed.} \(n \geq 5\) независимых запусков
    с фиксированными seed, отчётность по среднему и доверительному интервалу
    95\% (bootstrap, \(B = 1000\)).
    \item \textbf{Парные сравнения.} При сравнении пар моделей по seed
    применяется парный критерий Уилкоксона.
    \item \textbf{Запрет утечки.} Train/val/test-разделение по времени;
    препроцессорные статистики --- только по train.
    \item \textbf{Калибровка.} Все модели с поддержкой
    \(\hat{p}\) проходят температурное масштабирование на val-выборке;
    отчётность по ECE до и после.
    \item \textbf{Воспроизводимость.} Версионирование кода (commit hash),
    конфигов и данных; журнал \texttt{run\_id} для каждого запуска.
\end{enumerate}

\subsection{Базовые модели для сравнения}

Базовые модели (\emph{baselines}), используемые в эксперименте E1, и их
роль в работе приведены в \tabref{tab:baselines}.

\begin{table}[H]
\centering
\caption{Базовые модели сравнения, используемые в эксперименте E1}
\label{tab:baselines}
\renewcommand{\arraystretch}{1.25}
\begin{tabularx}{\linewidth}{@{}lXX@{}}
\toprule
\textbf{Модель} & \textbf{Тип} & \textbf{Назначение} \\
\midrule
Logistic Regression &
линейная supervised & нижняя граница линейного классификатора \\
Random Forest~\cite{breiman2001rf} &
ансамбль supervised & сильный baseline на табличных данных \\
XGBoost~\cite{chen2016xgboost} &
бустинг supervised & сильный baseline на табличных данных \\
MLP &
DL supervised & DL-baseline без учёта структуры \\
1D-CNN &
DL supervised & DL-baseline с локальной структурой \\
LSTM~\cite{hochreiter1997lstm} &
DL supervised & DL-baseline с временной структурой \\
GRU~\cite{cho2014gru} &
DL supervised & облегчённая RNN \\
BiLSTM &
DL supervised & альтернативная сильная DL-архитектура \\
TCN~\cite{bai2018tcn} &
DL supervised & альтернатива RNN \\
Transformer-encoder~\cite{vaswani2017attention} &
DL supervised & attention-based альтернатива \\
Autoencoder~\cite{mirsky2018kitsune} &
DL semi-supervised & контрольный AE \\
VAE~\cite{kingma2014vae, lopez2017cvae} &
DL semi-supervised & контрольный VAE \\
Isolation Forest~\cite{liu2008isolation} &
unsupervised & классический unsupervised \\
One-Class SVM~\cite{scholkopf2001ocsvm} &
unsupervised / one-class & классический one-class \\
\bottomrule
\end{tabularx}
\end{table}

\subsection{Сводная диаграмма экспериментального стенда}

\begin{figure}[H]
\centering
\begin{minipage}{0.94\linewidth}
\small
\begin{verbatim}
+----------------------------+        +----------------------------+
|  UNSW-NB15  (raw CSV)      |        |  CICIDS2017 (raw CSV)      |
+--------------+-------------+        +--------------+-------------+
               | preprocess + windowing              | preprocess + map
               v                                     v
+----------------------------+        +----------------------------+
|  Train / Val / Test         |       |  CICIDS2017 test (cross)   |
+--------------+--------------+       +--------------+--------------+
               | E1, E2, E3                          | E4'
               v                                     v
+--------------------------------------------------------------+
|              Models pool (Tab. 3.4): CNN-LSTM, ...           |
+----+--------------------+-----------------+------------+-----+
     | E1                 | E3              | E5        | E6
     v                    v                 v           v
+----------+        +----------+      +-----------+   +-----------+
| Compare  |        | Calibr.  |      | Drift inj |   | Lat./Thr. |
+----------+        +----------+      +-----------+   +-----------+
     ^                   ^                  ^             ^
     | E4                                                 |
     |                                                    |
+----+-----------------------------+                      |
|  AI Attacker (templates+FSM+drift) -------- stream ----+
+----------------------------------+
\end{verbatim}
\end{minipage}
\caption{Концептуальная схема экспериментального стенда}
\label{fig:experiment-stand}
\end{figure}

\subsection{Критерии успеха Этапа~4}

Этап~4 считается выполненным, если выполнены следующие условия.

\begin{enumerate}
    \item Полностью проведены эксперименты E1--E6, оформлены таблицы и
    графики метрик, графики обучения, confusion matrix.
    \item Гипотезы $H_0^{\text{main}}$, $H_0^{\text{aux}}$ и
    $H_0^{\text{transfer}}$ либо подтверждены при выполнении соответствующих
    NFR, либо обоснованно опровергнуты с указанием конкретных условий
    невыполнения.
    \item Сформирован итоговый отчёт с разделом ``Ограничения и направления
    дальнейшего развития'', явно фиксирующим невыполненные пункты и причины.
\end{enumerate}

\subsection{Промежуточные выводы по разделу}

\begin{enumerate}
    \item Сформирована библиотека из семи параметризованных шаблонов
    атакующих сценариев, согласованных с таксономией UNSW-NB15, с явной
    спецификацией параметров и привязкой каждого шаблона к классу датасета.
    \item Описаны конечно-автоматная политика агента, протокол валидации
    генератора (маргинальное соответствие нормы по KS и MMD;
    семантическая согласованность атак по средним и медианам).
    \item Сформирован протокол шести экспериментов (E1--E6 + E4') с
    явными целями, метриками и критериями принятия. Для всех экспериментов
    задана статистическая методология (множественные seed, bootstrap CI,
    парный критерий Уилкоксона).
    \item Зафиксирован пул из 14 моделей-кандидатов и baselines с
    указанием их роли в эксперименте E1; все модели обучаются и
    оцениваются в равных условиях, что устраняет ограничения, выявленные
    в опубликованной NIDS-литературе.
    \item Приведена диаграмма экспериментального стенда, отражающая
    взаимодействие подсистем: данные (UNSW-NB15, CICIDS2017),
    модели, активный тестовый источник (ИИ-злоумышленник) и блоки
    калибровки, drift-инжекции и измерения производительности.
\end{enumerate}
